\section{Reasoning model and operational contract}
\label{sec:semantics}

Nibli accepts a formal predicate-call language, nibli KR, and compiles it to an
indexed first-order representation. Its executable semantics is a supported,
bounded reasoning procedure over that representation, rather than an unrestricted
first-order validity decision procedure. This distinction separates three objects:
the statement a person intends, the formula accepted by the compiler, and the
answer obtained by evaluating that formula against the current premises.
Deterministic compilation addresses the second object. It does not establish
agreement with the first.

\subsection{A finite comparison fragment}

We define the fragment used in the experiments independently of the implementation.
Let $C$ be the finite set of opaque constants occurring in a workload's program
and queries, and $\Sigma$ a finite set of predicate symbols of fixed arity.
A ground atom is $p(c_1,\ldots,c_k)$ with $c_i\in C$.
A program consists of a finite extensional set $E$ and range-restricted rules
\begin{equation}
 H \leftarrow A_1,\ldots,A_m,\ \mathop{\mathrm{not}} B_1,\ldots,
 \mathop{\mathrm{not}} B_n .
 \label{eq:rule}
\end{equation}
Every variable occurs in a positive body atom. There are no function symbols,
existential heads, arithmetic operations, equality merges, or modal operators in
this fragment. Extensional and intensional occurrences of the same predicate
are permitted, as in the reachability experiment.

For each rule, put an edge from its head predicate to every body predicate, with
negative polarity for a negated body atom. A stratification assigns a natural
number $s(p)$ to each predicate such that $s(H)\geq s(A_i)$ and $s(H)>s(B_j)$.
Nibli rejects registration that creates a strongly connected component containing
a negative edge. A rejection is an input error, not an \Unknown\ answer.
For accepted fragment programs, compute a least fixed point within each stratum
after fixing the extensions of lower strata. We write the resulting perfect
model as $M(E,P)$. The chain and policy workloads have hand-derived expected
memberships in this model; clingo and \Souffle\ supply independently executed
comparison results.

A negative body literal holds only when its positive atom is absent from the
completed lower-stratum model. This is the closed-world assumption. It is not a
claim that the physical world contains no further evidence. Similarly, queries
quantify over the engine's represented domain rather than every possible entity
in the world. Range restriction means that a fresh query-only constant cannot
create positive support merely by appearing in an adapter's symbol universe.
The comparison uses ground membership queries, not cross-system comparisons
of quantified domains. Shuffling input order or consistently renaming
those constants changes an execution workload without changing the expected
abstract membership answers.

\subsection{The broader KR surface}

The implemented surface includes quantified arguments, connected clauses,
equality, generated witnesses, tense, deontic wrappers, and opaque quoted
abstractions. These features make the integrated engine more than a direct
Datalog evaluator. They also prevent a claim that the entire surface has the
same semantics as the comparison fragment. The artifact's claim register
separates implementation tests for these features from the narrower external
solver comparison.

An ordinary multi-place predicate is represented through an event anchor and
role predicates. Schematically, a ground surface assertion $p(a,b)$ introduces
\begin{equation}
 \exists e.\ p(e)\land p_{x1}(e,a)\land p_{x2}(e,b).
 \label{eq:event}
\end{equation}
The actual representation carries typed terms and shared node indices. Equation
\ref{eq:event} is a readable projection, not a literal serialization of the
intermediate buffer. Rules can introduce dependent event witnesses. Consequently,
a function-free surface workload need not have a function-free internal encoding.
Both materialization eligibility and external solver translation must account
for this seam. Unspecified corpus places are also implicit in the KR compilation;
the comparison concerns the explicitly supplied argument positions, used
consistently in facts, rules, and queries.

Names, generated individuals, generated events, and compiler placeholders are
different term cases. A generated witness has a semantic identity containing its
source assertion ID, binder ordinal, sort, and origin; a dependent witness also
depends on its arguments. A printed label such as \texttt{sk\_0} is a display
convenience. It is not an address that an input string can use to select the
witness. This matters when a source is withdrawn or replayed: confusing a user
constant with a generated witness would create joins and citations that have
no supporting premise.

The represented domain includes named entities and admitted generated individuals,
with finite numeric terms occurring in assertions also contributing members.
Querying does not itself assert a new domain member. Equality can identify
represented entities, and withdrawal reconstructs the domain and equality indexes.
The default universal reading has no existential import: a universally stated
restriction does not by itself establish a member of its restrictor class.
An explicit legacy profile changes that behavior and is recorded in the proof
envelope. All new measurements use existential import disabled.

Tense wrappers distinguish stored and derived flavors; they do not imply a
clock, interval arithmetic, or automatic temporal lifting of bare rules. Opaque
abstractions describe content without asserting that the content actually holds.
These are representation commitments enforced by compilation and reasoning
tests. They are not claims to cover arbitrary natural-language tense or modal
logic. In particular, the finite comparison deliberately excludes those forms
rather than silently flattening them into ordinary predicates.

\subsection{Four outcomes, not four truth values}
\label{sec:outcomes}

The public result type has four operational classes, with reasons on the two
non-definitive classes (Table~\ref{tab:outcomes-contract}). Calling this an
operational interface avoids assigning an epistemic probability or a
many-valued model-theoretic meaning to a resource failure.

\begin{table}[t]
\centering
\small
\begin{tabular}{lp{10.5cm}}
\toprule
Outcome & Meaning for a supported query in its recorded profile \\
\midrule
\True & Evaluation finds a positive derivation, or an accepted trusted
compute result establishes the query. \\
\False & Evaluation definitively fails to derive the claim in the current
closed model, or a supported computation definitively refutes its supplied
arguments. It is not, in general, a derivation of explicit negation. \\
\Unknown & Evaluation cannot settle the claim for a recorded reason, such as
a cycle cut, unavailable compute backend, or non-finite computation. \\
\Resource & A configured depth or host resource limit prevents a definitive
answer. The reason distinguishes depth, fuel, or memory. \\
\bottomrule
\end{tabular}
\caption{Operational result contract. An external experiment cutoff is recorded
outside this table; it is not an engine result.}
\label{tab:outcomes-contract}
\end{table}

Negation must preserve this distinction. If proving $A$ remains non-definitive,
its absence has not been established, so a negation-as-failure check cannot
turn that uncertainty into a definitive success. Connective combination also
has decisive cases: a false conjunct settles a conjunction, while a true
disjunct settles a disjunction. For unresolved combinations, resource exhaustion
takes precedence over uncertainty. Reasons need not be commutative even when
the result class is. These finite combinations are the part of the Lean-to-Rust
correspondence that can be checked exhaustively.

Depth is an evaluation horizon, not a limit on the number of edges appearing
in a dataset. A transitive rule can compose two already established subpaths.
Materialization can also establish a complete finite extension that settles a
query beyond the backward search horizon. Thus a chain longer than the configured
depth need not return \Resource, and merely increasing the depth does not
characterize elapsed time. The experiments measure definitive coverage alongside
latency to expose this distinction.

\subsection{Premise withdrawal and the trusted boundary}
\label{sec:withdrawal}

Let $R$ be an ordered registry of accepted assertion records with stable IDs, and
let $L(R)$ be its active subsequence. Ordinary withdrawal changes a record's
status and reconstructs logical state by replaying $L(R)$, preserving record
identity. For a query $q$, the intended observational property is
\begin{equation}
 \operatorname{eval}(\operatorname{withdraw}(R,i),q,\pi)
 \simeq
 \operatorname{eval}(\operatorname{replay}(L(R)\setminus\{i\}),q,\pi),
 \label{eq:replay}
\end{equation}
where $\pi$ fixes the session profile. Verdict equivalence is the central
observation; stable source IDs permit the evidence to be checked against the
surviving registry as well. This property does not say that withdrawing one
support removes every derivation. Alternative support must remain effective.

Some declarations intentionally survive rebuild, including restrictions on
which predicates may be asserted directly. Equation~\ref{eq:replay} therefore
applies to ordinary premises under the same admission state and configuration,
not to treating every administrative declaration as a reversible logical fact.
Reset and reopen are separate operations. Reopen reconstructs active logical
records and retains withdrawn IDs; it is not a replay of a complete chronological
audit log. The new update experiment explicitly supplies its operation sequence
because retained status alone cannot recover that history.

The compute interface is another boundary of the argument. A backend evaluates
supplied arguments and its reply is trusted for the current query. The reply is
not automatically inserted into the durable premise registry. Failure to obtain
a supported reply remains non-definitive. A computation step in a trace is
evidence that the engine used that boundary, not a proof that the external
implementation, input source, or freshness policy is correct. The current
envelope does not supply a cryptographically authenticated backend transcript.

These qualifications give a precise interpretation of the project's
\enquote{zero hallucination} description: conclusions are constrained by accepted
premises, supported inference, the recorded evaluation profile, and explicitly
trusted computation. The phrase does not certify factual inputs, intended
natural-language meaning, every implementation path, or unrestricted logical
completeness. A wrong premise can yield a correctly derived wrong real-world
decision.
